\chapter{Invariance causale des interfaces par rapport aux références externes}
\label{chap:independencia-prospectiva}
\index{indépendance prospective}\index{graphe causal}
\index{protocolo ciego}\index{tours!complétude reconstructive}

\section{Critère d'indépendance causale par rapport aux références externes : graphe dirigé acyclique et fixation préalable de la sortie}
\label{sec:definiciones-independencia-prospectiva}

Soient \(\mathsf P\) les primitives internes, \(\mathsf R\) les règles,
lecteurs, calibres et normalisations déclarés, \(G\) un générateur,
\(o=G(\mathsf P,\mathsf R)\) sa sortie et \(y\) une référence externe.

\begin{definition}[Interface prospectivement indépendante]
\label{def:interfaz-prospectivamente-independiente}
L'interface est prospectivement indépendante de \(y\) si :
\begin{enumerate}
 \item \(y\notin\operatorname{Anc}(o)\);
 \item \(\mathsf P,\mathsf R,G\) et les contrôles négatifs sont figés avant
 l'ouverture de \(y\) ;
 \item l'exécution n'a accès ni directement ni indirectement à \(y\) ;
 \item \(o\) et sa trace sont scellés avant le calcul d'une comparaison
 \(\operatorname{cmp}(o,y)\).
\end{enumerate}
\end{definition}

\begin{definition}[Indépendance de la sélection]
\label{def:independencia-seleccion}
Le caractère aveugle de chaque exécution ne suffit pas si, après avoir produit des candidats
\(G_j(\mathsf P,\mathsf R_j)\), on publie
\[
 j_\ast(y)=\arg\min_j
 \operatorname{cmp}\!\left(G_j(\mathsf P,\mathsf R_j),y\right).
\]
L'indépendance de la sélection interdit que cette opération détermine la sortie
publiée.
\end{definition}

\begin{definition}[Prédiction multivaluée]
\label{def:prediccion-multivaluada}
Si les descripteurs physiques non métrologiques et les restrictions internes
admettent plusieurs multisections, la sortie prospective est l'ensemble complet
\[
 \widehat{\mathfrak R}(d)
 =\{\mathcal M:\mathcal M\text{ satisfait toutes les contraintes
 préinscrites pour }d\}.
\]
La non-unicité est une sortie de la procédure. Choisir ensuite un élément
par proximité d'une grandeur est interdit.
\end{definition}

\section{Théorème de non-interférence}
\label{sec:teorema-no-interferencia}

\begin{definition}[Graphe causal dirigé acyclique fidèle et complet]
\label{def:dag-fiel-completo}
Un graphe dirigé acyclique \(D\) est fidèle si chaque nœud représente la même
fonction et les mêmes entrées que la construction citée. Il est complet pour
une sortie si toute dépendance causalement pertinente apparaît comme arête ou
racine déclarée. La fidélité et la complétude sont des hypothèses sémantiques ;
l'acyclicité de l'ensemble de données ne les démontre pas.
\end{definition}

\begin{theorem}[Indépendance prospective conditionnée par le graphe causal]
\label{thm:no-interferencia-prospectiva}
Soit un système déterministe d'équations structurelles représenté par un
graphe dirigé acyclique fini \(D\), fidèle et complet pour \(o\). Si
\(y\notin\operatorname{Anc}_D(o)\), toutes les racines ancestrales de \(o\)
restent figées et l'exécution n'a pas d'accès caché à \(y\), alors
pour tous \(y_0,y_1\),
\[
 o_{\operatorname{do}(y=y_0)}
 =o_{\operatorname{do}(y=y_1)}.
\]
\end{theorem}

\begin{proof}
On ordonne topologiquement les antécédents de \(o\). Leurs racines sont internes
et restent fixes. Par induction sur l'ordre, chaque nœud reçoit les mêmes
arguments sous les deux interventions et, par déterminisme, renvoie la même
valeur. L'égalité se propage jusqu'à \(o\). L'empreinte cryptographique ne crée pas
l'indépendance ; elle rend seulement détectable une modification ultérieure.
\end{proof}

\begin{corollary}[Permutation de la comparaison externe]
\label{cor:permutacion-comparacion-externa}
Si les noms, les lignes ou les valeurs externes sont permutés après le scellement, la
sortie interne ne change pas, bien que son score puisse changer :
\[
 o_{\operatorname{do}(y=\sigma y_0)}
 =o_{\operatorname{do}(y=y_0)},\qquad
 \operatorname{cmp}(o,\sigma y_0)
 \ne\operatorname{cmp}(o,y_0).
\]
\end{corollary}

\begin{proposition}[Séparation des branches calibrées]
\label{prop:separacion-ramas-calibradas}
Si une référence externe est intervenue dans une sélection historique, la
dépendance ne contamine pas les objets internes qui ne descendent pas de cette
sélection. La branche peut être isolée et soumise à un nouveau protocole prospectif
sans modifier le noyau.
\end{proposition}

\begin{proof}
Être un antécédent est une propriété locale du graphe. Les branches
\[
 Z_0\longrightarrow p=5\longrightarrow\text{échelle électronique},
\qquad
 \{\text{masses}\}\longrightarrow\{\text{signatures calibrées}\}
 \longrightarrow\text{table}
\]
restent hors du sous-graphe protégé. Les tours, le caractère formel,
l'atlas et le centre \((5,5)\) ne reçoivent pas d'arêtes de ces références.
\end{proof}

\section{Invariance contrefactuelle sur le graphe causal dirigé et portée de la conclusion}
\label{sec:certificado-independencia-prospectiva}

Le graphe causal publié contient \(71\) nœuds, \(89\) arêtes,
\(16\) interfaces et
\(29\) sorties protégées. La vérification causale établit l'acyclicité, les décomptes,
les voies déclarées, l'existence des sources enregistrées, l'absence de
références externes parmi les antécédents consignés et l'invariance
contrefactuelle du modèle symbolique. Le résultat est
\[
 \#\{\text{prédécesseurs externes de sorties protégées}\}=0.
\]
Le certificat ne démontre pas à lui seul qu'aucune arête ne manque et n'audite pas
toutes les lectures de données de chaque programme scientifique. Sa force est donc
\textsc{démontré avec une structure de départ explicite} : la conclusion
s'obtient sur le graphe dirigé acyclique expressément déclaré et reste
conditionnée par sa fidélité et sa complétude.

Le remplacement simultané des références CODATA, Planck,
Barbero--Immirzi, CKM, masses, gravité et cristal temporel peut modifier les
comparaisons et les branches calibrées. Il ne modifie pas, dans le graphe déclaré,
le catalogue de \(468\) émissions, les cinq régions de \(\pi\), la
génération coinductive corrélée, \(\alpha_{12}\), \(C^\ast\), la décade \(-34\), les tours,
l'atlas, le centre électronique, la double projection ni le corridor
exceptionnel.

\section{Essai multivalué sans accès à la valeur cible et critère d'identification}
\label{sec:protocolo-ciego-multivaluado}

\begin{algorithm}[Attribution future sans grandeurs]
\label{alg:asignacion-futura-sin-magnitudes}
L'interface physique s'exécute en quatre phases irréversibles :
\begin{enumerate}
 \item \emph{Gel} : on scelle l'autorité, l'atlas, les involutions,
 les restrictions, le code, la priorité et les contrôles négatifs.
 \item \emph{Entrée autorisée} : seuls sont reçus le secteur, la génération, la charge,
 la couleur, la conjugaison, le spin/la statistique, la composition et le type métrologique ;
 les masses, les signatures raffinées, les résidus, les angles et les unités sont exclus.
 \item \emph{Sortie} : toutes les multisections compatibles sont énumérées ; pour
 chaque extension de chaque multisection, on calcule la voie observable, le registre,
 la signature, \(\chi_{\rm form}\), la spécialisation et l'observable. Dans les objets
 composés, on raccorde d'abord les extensions et les registres. On scelle le
 multiensemble complet.
 \item \emph{Ouverture} : un dépositaire distinct ouvre les grandeurs et calcule
 les résidus, les incertitudes et les covariances sans modifier aucun candidat.
\end{enumerate}
\end{algorithm}

La multisection enrichie, et non une extension ponctuelle sélectionnée
rétrospectivement, est la sortie de l'appariement physique. L'évaluation
ultérieure peut être effectuée fibre par fibre. L'algorithme précédent est une
spécification réfutable ; il n'est pas présenté comme une prédiction multiensemble déjà
exécutée.

\begin{criterion}[Échec du protocole en aveugle]
\label{crit:fallo-protocolo-ciego}
Le protocole échoue si le code peut lire une grandeur, si une signature
historique calibrée réapparaît comme entrée, si un candidat est écarté après
ouverture de la référence, si une masse positive est confondue avec une masse exactement nulle
ou si le multiensemble scellé change lorsque le catalogue externe est permuté.
\end{criterion}

\section{Caractère formel et complétude reconstructive}
\label{sec:completitud-reconstructiva-torres}

Pour \(v=(n_A,s_C,k,b_{120},b_{\zeta})\in\ZZ^5\), le caractère formel
antérieur à toute spécialisation est
\[
 \chi_{\rm form}(v;x_A,x_C,x_\Delta,x_{120},x_{270})
 =\exp\!\left(
 n_Ax_A+s_Cx_C+kx_\Delta
 +b_{120}x_{120}+b_{\zeta}x_{270}\right).
\]
La spécialisation déclarée
\[
 \Theta_0
 =\left(A_{\rm rad},\frac{C^\ast_{\rm rad}}6,
 \Dfour,s_{120},\zeta_{270}\right),
 \qquad \zeta_{270}=135\Delta_4^2,
\]
produit
\(\chi_0(v)
=\chi_{\rm form}(v;\Theta_0)\).
La signature brute, le caractère formel, sa spécialisation centrale, le
raffinement global \(\eta\in\mathcal E_{\mathfrak S}\) et un tableau
calibré sont cinq objets différents.

\begin{theorem}[Complétude reconstructive positive]
\label{thm:completitud-reconstructiva-positiva}
Soit
\[
 \mathcal N=\langle90,120\rangle_{>0}
\]
et soit \((n_j)_{j\ge1}\) son énumération croissante. Étant donnés
\(0<q_+<q_-<1\), définissons
\[
 a_j=q_-^{n_j}-q_+^{n_j}>0.
\]
Alors \(a_j\to0\) et \(\sum_ja_j<\infty\). Une règle déterministe étant fixée
pour départager les demi-entiers, pour tout \(r\in\RR\), la récurrence
\[
 r_0=r,\qquad
 \nu_j=\operatorname{nint}\!\left(\frac{r_{j-1}}{a_j}\right),\qquad
 r_j=r_{j-1}-\nu_ja_j
\]
satisfait
\[
 |r_j|\leq\frac{a_j}{2}\longrightarrow0,\qquad
 r=\sum_{j\ge1}\nu_ja_j,\qquad
 \sum_{j\ge1}|\nu_ja_j|<\infty.
\]
Ainsi, tout \(R>0\), avec \(r=\log R\), est représenté exactement dans la
complétion par
\(\exp(\sum_j\nu_ja_j)\).
\end{theorem}

\begin{proof}
Les indices de \(\mathcal N\) sont \(90,120\) et tous les multiples de \(30\)
à partir de \(180\). Par conséquent,
\[
 \sum_{n\in\mathcal N}q_-^n
 =q_-^{90}+q_-^{120}
 +\frac{q_-^{180}}{1-q_-^{30}}<\infty.
\]
Comme \(0<a_j<q_-^{n_j}\), la suite est sommable et converge vers zéro. La
définition de l'entier le plus proche implique
\(|r_j|\le a_j/2\). Par télescopage,
\(\sum_{i=1}^{N}\nu_ia_i=r-r_N\to r\). Pour \(j\ge2\),
\[
 |\nu_ja_j|=|r_{j-1}-r_j|
 \le\frac{a_{j-1}+a_j}{2},
\]
et la sommabilité de \((a_j)\) prouve la convergence absolue.
\end{proof}

Le théorème est une complétude du dictionnaire sur \(\RR_{>0}\). Si les
coefficients \(\nu_j\) sont obtenus à partir du résidu d'une masse connue, leur utilisation
est une \textsc{reconstruction calibrée} ; elle ne fournit pas la sélection
prospective. Elle ne représente pas non plus une masse exactement nulle. Le photon, le gluon ou un
éventuel mode de spin deux exigent une protection exacte par symétrie dans la
représentation physique.

\Needspace{34\baselineskip}
\section{Protocoles prospectifs d'invariance pour les familles de constantes}
\label{sec:protocolos-prospectivos-especificos}

\begin{description}
 \item[Barbero--Immirzi.] On fige \(S_{90},S_{120}\), les poids
 \((12,-1)\) et leur normalisation avant d'ouvrir une référence physique.
 L'évaluation de \(\Gamma_{6\to8}\) est aveugle à cette référence,
 mais l'unicité préalable des séries et des poids et l'opérateur d'entrelacement avec l'opérateur
 d'aire sont des problèmes distincts.
 \item[CKM et CP.] On scelle les quatre équations, leurs coefficients et la
 convention angulaire ; on publie simultanément les quatre paramètres,
 \(|V|\) et \(J\). La confrontation utilise une covariance fixée à l'avance.
 \item[Constitutives.] Les identités adimensionnelles précèdent
 \(Z_{\rm ref}\) et \(c_{\rm ref}\). Le changement de carte transforme la
 réalisation, non les identités, et ne crée pas quatre prédictions
 indépendantes.
 \item[Gravité et inertie.] Les rapports
 \(G_{\rm tor}/I_{\rm tor}=1\) et
 \(G_{\rm av}/I_{\rm av}=\pi/\varphi^2\) appartiennent à des couches différentes.
 Leur interprétation physique exige un morphisme et une confrontation prospective ; un
 phénomène cosmologique ne sélectionne pas le rapport.
\end{description}

\section{Stratification probatoire : identités internes, déductions, reconstructions calibrées et interfaces soumises à une fermeture mathématique}
\label{sec:cinco-niveles-independencia}

\begin{table}[htbp]
\centering
\caption{Niveaux que l'indépendance causale ne permet pas de confondre.}
\label{tab:cinco-niveles-independencia}
\small
\begin{tabularx}{\textwidth}{@{}L{0.25\textwidth}YL{0.27\textwidth}@{}}
\toprule
Niveau & Exemple & Énoncé autorisé\\
\midrule
Générateur interne certifié&
Régions, \(\alpha_{12}\), \(C^\ast\), \(-34\), atlas&
Sortie obtenue dans le domaine déclaré.\\
Évaluation conditionnée en aveugle&
\(\Gamma_{6\to8}\), CKM fixé, \(\chi_0(v)\)&
Ne consulte pas la cible, mais dépend de la fonctionnelle ou de la signature.\\
Protocole en aveugle non exécuté&
Spécification multivaluée&
Il existe une procédure réfutable ; non une sortie scellée.\\
Reconstruction calibrée&
Signatures historiques et raffinement à partir de \(r=\log R\)&
Évaluation exacte conditionnée par une référence.\\
Interface physique en attente&
Particule--antiparticule, aire, masse nulle, gravité&
Domaine défini ; représentation ou confrontation manquante.\\
\bottomrule
\end{tabularx}
\end{table}
